% ============================================================
% f2013_zhejiang_science_q20_tetrahedron.tex
% 2013 浙江卷（理科）第 20 题 四面体 A-BCD（TikZ 矢量重绘）
% 题设：AD⊥平面 BCD，BC⊥CD，AD=2，BD=2√2；
%       M 为 AD 中点，P 为 BM 中点，Q 在 AC 上且 AQ=3QC。
% 实线：可见棱 AB、AC、AD、BC、CD 及辅助线 CM；
% 虚线：被遮挡的 BD、BM（过 P）、PQ。
% ============================================================
\begin{tikzpicture}[
  x=1cm, y=1cm,
  every node/.style={font=\small, inner sep=1pt},
  s/.style={line width=0.6pt},
  d/.style={line width=0.6pt, dash pattern=on 3.0pt off 2.4pt},
]

  \coordinate (A) at (5.70,6.20);
  \coordinate (B) at (0.20,1.60);
  \coordinate (C) at (3.50,0.15);
  \coordinate (D) at (5.90,1.90);

  % ---------- 实线 ----------
  \draw[s] (A) -- (B);          % AB
  \draw[s] (A) -- (C);          % AC（过 Q）
  \draw[s] (A) -- (D);          % AD（过 M）
  \draw[s] (B) -- (C);          % BC
  \draw[s] (C) -- (D);          % CD

  % ---------- 辅助线与虚线 ----------
  \draw[s] (C) -- ($(A)!0.5!(D)$);  % CM（面 ACD 内的辅助线）
  \draw[d] (B) -- (D);              % BD（被遮挡）
  \draw[d] (B) -- ($(A)!0.5!(D)$);  % BM（过 P）

  % ---------- 点 ----------
  \coordinate (M) at ($(A)!0.5!(D)$);
  \coordinate (P) at ($(B)!0.5!(M)$);
  \coordinate (Q) at ($(A)!0.75!(C)$);
  \fill (M) circle [radius=0.035];
  \fill (P) circle [radius=0.035];
  \fill (Q) circle [radius=0.035];
  \draw[d] (P) -- (Q);              % PQ（体内线段）

  % ---------- 标签 ----------
  \node[above] at (A) {$A$};
  \node[left] at (B) {$B$};
  \node[below] at (C) {$C$};
  \node[right] at (D) {$D$};
  \node[right] at (M) {$M$};
  \node[above left] at (P) {$P$};
  \node at ($(Q)+(0.26,-0.30)$) {$Q$};

\end{tikzpicture}
